\documentclass{amsart}
\usepackage{amsmath,amssymb,amsthm,hyperref}
\usepackage{algorithm}
\usepackage{algpseudocode}
\newtheorem{theorem}{Theorem}
\newtheorem{lemma}{Lemma}
\newtheorem{proposition}{Proposition}
\theoremstyle{definition}
\newtheorem{definition}{Definition}
\newtheorem{remark}{Remark}

\begin{document}

\title{Worst-case complexity of collection in polycyclic groups}
\maketitle

\section{Statement of the result}

We work with a consistent polycyclic presentation of a polycyclic group $G$ of
Hirsch length $n$, given by generators $g_1,\dots,g_n$, relative orders
$r_i\in\{2,3,\dots\}\cup\{\infty\}$, power relations $g_i^{r_i}=w_{i,i}$ (for finite
$r_i$) and conjugation relations $g_j^{g_i}=w_{i,j}$, $g_j^{g_i^{-1}}=w'_{i,j}$ for
$1\le i<j\le n$, each right-hand side being a \emph{normal word}
$g_{i+1}^{e_{i+1}}\cdots g_n^{e_n}$ with $0\le e_k<r_k$ whenever $r_k<\infty$. We use
the convention that the normal form lists generators in \emph{increasing} index from
left to right,
\[
g_1^{a_1}g_2^{a_2}\cdots g_n^{a_n},\qquad 0\le a_i<r_i\ (r_i<\infty),\ a_i\in\mathbb Z\ (r_i=\infty).
\]

Throughout, the cost model is the bit (Turing-machine) model. Let
\[
b \;=\; \text{the maximum bit-length of any integer exponent}
\]
occurring in the relator words $w_{i,i},w_{i,j},w'_{i,j}$ \emph{or} in the input word,
and let $M=2^{b}$, so that every individual exponent in the relators and in the input has
absolute value $<M$. The input is a word
\[
W \;=\; g_{c_1}^{\,\varepsilon_1}\cdots g_{c_\ell}^{\,\varepsilon_\ell},
\qquad c_t\in\{1,\dots,n\},\ \varepsilon_t\in\mathbb Z\setminus\{0\},
\]
of \emph{syllable length} $\ell$ (the number of generator powers written down); the
bit-lengths of the input exponents $\varepsilon_t$ are absorbed into $b$.

\begin{theorem}\label{thm:main}
Let $G$ be given by a consistent polycyclic presentation of Hirsch length $n$ as above,
with relator/input exponent bit-bound $b$, and let $W$ be an input word of length $\ell$.
Then the collection-from-the-left algorithm \textup{(}Algorithm~\textup{\ref{alg:collect})}
computes the normal form of $W$, and:
\begin{enumerate}
\item \textup{(Exponent size.)} Every integer exponent that appears at any moment during the
computation has absolute value at most
\begin{equation}\label{eq:expbound}
B \;:=\; \ell\,M\,(1+nM)^{\,n-1},
\end{equation}
i.e.\ bit-length $\;\log_2 B=O\!\bigl(n\,(b+\log n)+\log\ell\bigr)=O\!\bigl(n\,(b+\log(n\ell))\bigr).$
\item \textup{(Termination and upper bound.)} The algorithm halts, and its running time is
bounded by
\begin{equation}\label{eq:upper}
\ell\cdot 2^{\,O\!\left(n^2(b+\log(n\ell))\right)}\;=\;(nM\ell)^{\,O(n^2)},
\end{equation}
which is polynomial in $\ell$ and in $M=2^{b}$ for each fixed Hirsch length $n$, and
exponential in $n$.
\item \textup{(Lower bound / tightness in $n$.)} The exponential dependence on $n$ is
unavoidable: there is a family of consistent polycyclic presentations \textup{(}all relator
exponents in $\{-1,0,1\}$, so $b=1$\textup{)} of Hirsch length $N=\binom{n+1}{2}$, and inputs
of length $\ell=n^2$, for which the normal form contains a single exponent equal to
$\binom{2n-1}{n}$, of bit-length $2n-O(\log n)=\Theta(\sqrt N)$. Hence merely writing the
output, and a fortiori the running time of the collection algorithm, is exponential in $n$ on
these instances although $b=1$ and $\ell=\mathrm{poly}(n)$.
\end{enumerate}
\end{theorem}

\noindent
Thus the precise answer to ``determine the worst-case complexity'' is
\[
\boxed{\;(nM\ell)^{\,O(n^2)}\;=\;\ell\cdot 2^{\,O(n^2(b+\log(n\ell)))}\;,\qquad M=2^{b},\;}
\]
\emph{polynomial in the word length $\ell$ and in the exponent magnitude $M=2^{b}$ for each
fixed Hirsch length, and exponential in the Hirsch length $n$; the exponential blow-up in $n$
is genuine.}

\medskip
We prove the theorem as follows: \S\ref{sec:alg} specifies the algorithm and its correctness;
\S\ref{sec:exp} proves the exponent bound \eqref{eq:expbound} (this needs no prior knowledge
of termination); \S\ref{sec:steps} bounds the number of elementary operations, which both
proves termination and yields \eqref{eq:upper}; \S\ref{sec:lower} gives the lower-bound
construction.

\section{The algorithm and its correctness}\label{sec:alg}

We represent a group word as a finite list of \emph{syllables}, each syllable a pair $(c,e)$
standing for $g_c^{\,e}$ with $c\in\{1,\dots,n\}$ and $e\in\mathbb Z\setminus\{0\}$. The
\emph{level} of $(c,e)$ is its generator index $c$. A list $(c_1,e_1),\dots,(c_m,e_m)$ is
\emph{collected} (in normal form) if $c_1<c_2<\cdots<c_m$ and, for each $t$, $0\le
e_t<r_{c_t}$ whenever $r_{c_t}<\infty$.

The algorithm keeps a \emph{collected prefix} $P$ (always in normal form, with maximal level
$\mathrm{top}(P)$, taken $0$ when $P$ is empty) and an \emph{uncollected tail} $U$ (a list),
and repeatedly processes the first syllable $(c,e)$ of $U$.

\begin{algorithm}[h]
\caption{Collection from the left}\label{alg:collect}
\begin{algorithmic}[1]
\Require A consistent polycyclic presentation $(g_i,r_i,w_{i,i},w_{i,j},w'_{i,j})$ and a
word $W$ as a syllable list.
\Ensure The normal form of $W$.
\State $P\gets$ empty list;\quad $U\gets W$.
\While{$U$ is nonempty}
  \State let $(c,e)$ be the first syllable of $U$; remove it from $U$.
  \If{$e=0$} \State \textbf{continue} \EndIf
  \If{$c\le\mathrm{top}(P)$}\label{ln:swap}
     \State let $(k,f)$ be the last syllable of $P$, so $k=\mathrm{top}(P)\ge c$; remove it
     from $P$.
     \If{$c=k$} \State prepend $(c,f+e)$ to $U$ \Comment{merge}
     \Else \Comment{$c<k$: move one $g_c$ leftwards}
        \State let $\sigma\in\{+1,-1\}$ be the sign of $e$;
        \State prepend to $U$, in order: $(c,\sigma)$, then the word $\mathrm{Conj}(g_k^{f},g_c^{\sigma})$,
        then $(c,e-\sigma)$.\label{ln:swapprepend}
     \EndIf
  \ElsIf{$r_c=\infty$ \textbf{or} $0\le e<r_c$}\label{ln:place}
     \State append $(c,e)$ to the end of $P$ \Comment{$(c,e)$ is now correctly placed}
  \Else \Comment{$c>\mathrm{top}(P)$ but $e$ out of range; reduce the power}\label{ln:power}
     \State write $e=q\,r_c+e_0$ with $0\le e_0<r_c$ (so $q\ne0$);
     \State prepend to $U$, in order: $(c,e_0)$, then $q$ copies of $w_{c,c}$ if $q>0$, or
     $|q|$ copies of $w_{c,c}^{-1}$ if $q<0$.\label{ln:powerprepend}
  \EndIf
\EndWhile
\State \Return $P$.
\end{algorithmic}
\end{algorithm}

Here $\mathrm{Conj}(g_k^{f},g_c^{\sigma})$ denotes the word obtained from $g_c^{-\sigma}\,g_k^{f}\,g_c^{\sigma}$
by substituting, in each of the $|f|$ letters $g_k^{\pm1}$, the relator
$g_k^{\,g_c^{\sigma}}$; explicitly it is the concatenation of $f$ copies of $w_{c,k}$ (if
$\sigma=+1$, $f>0$), of $|f|$ copies of $w_{c,k}^{-1}$ (if $\sigma=+1$, $f<0$), or with
$w_{c,k}$ replaced by $w'_{c,k}$ when $\sigma=-1$. By the defining property of a polycyclic
presentation, every $w_{c,k},w'_{c,k}$ is a normal word in $g_{c+1},\dots,g_n$; likewise
$w_{c,c}$ is a normal word in $g_{c+1},\dots,g_n$. The formal inverse $w^{-1}$ of a word $w$
is the reversed word with each syllable exponent negated. We record the structural fact used
throughout.

\begin{lemma}[Level monotonicity]\label{lem:level}
Consider any single iteration of Algorithm~\ref{alg:collect} acting on a level-$c$ syllable.
Then:
\begin{enumerate}
\item All syllables in the words $\mathrm{Conj}(g_k^{f},g_c^{\sigma})$ \textup{(}swap with
$c<k$\textup{)} and in the copies of $w_{c,c}^{\pm1}$ \textup{(}power reduction\textup{)} have
level $\ge c+1$.
\item Apart from those words, the only syllables created have level $c$ \textup{(}namely the
leading $(c,\sigma)$ and residual $(c,e-\sigma)$ in a swap, or the residual $(c,e_0)$ in a
power reduction, or the merged $(c,f+e)$\textup{)}.
\end{enumerate}
In particular, no iteration acting on a level-$\ge i$ syllable ever creates a syllable of
level $<i$.
\end{lemma}

\begin{proof}
Immediate from lines~\ref{ln:swapprepend} and \ref{ln:powerprepend} and the definition of
$\mathrm{Conj}$, since every $w_{c,k},w'_{c,k},w_{c,c}$ is a normal word in
$g_{c+1},\dots,g_n$. The last sentence follows because an iteration on a level-$i$ syllable is
a placement (creates nothing), a power reduction at level $i$ (creates levels $\ge i+1$ and a
level-$i$ residual), a merge at level $i$ ($k=i$, creates one level-$i$ syllable), or a swap
with $c=i<k$ (creates levels $\ge i+1$ and two level-$i$ syllables); none produces level $<i$.
\end{proof}

\paragraph{Correctness.}
Each iteration rewrites $P\cdot U$ to an equal element of $G$. Placement and the various
prepends merely re-bracket a word. The merge uses $g_k^{f}g_c^{e}=g_c^{f+e}$ when $c=k$. The
swap uses the identity $g_k^{f}g_c^{\sigma}=g_c^{\sigma}\,\mathrm{Conj}(g_k^{f},g_c^{\sigma})$,
which is exactly the conjugation relation $g_k^{\,g_c^{\sigma}}=w_{c,k}$ (resp.\ $w'_{c,k}$)
applied $|f|$ times; consequently $g_k^{f}g_c^{e}=(c,\sigma)\cdot
\mathrm{Conj}(g_k^{f},g_c^{\sigma})\cdot(c,e-\sigma)$ as group elements. The power reduction
uses $g_c^{r_c}=w_{c,c}$, hence $g_c^{e}=g_c^{e_0}(g_c^{r_c})^{q}=g_c^{e_0}w_{c,c}^{q}$. Thus
the element of $G$ represented by $P\cdot U$ is invariant. When the loop ends, $U$ is empty and
$P$ is collected, so $P$ is \emph{a} normal word equal to the input element. Consistency of the
presentation guarantees the normal word representing a given element is unique; hence $P$ is
\emph{the} normal form of $W$. It remains to bound exponents (\S\ref{sec:exp}) and prove the
loop terminates with the stated number of iterations (\S\ref{sec:steps}).

\section{Exponent size bound}\label{sec:exp}

For a configuration $(P,U)$ and a level $c\in\{1,\dots,n\}$ define the \emph{level-$c$ mass}
\[
N_c\;=\;\sum_{(c,e)\in P\cup U}|e|.
\]
The bound below is proved by causal induction on levels and does not assume termination.

\begin{proposition}\label{prop:exp}
Define $A_1=\ell M$ and $A_c=\ell M+nM\sum_{j<c}A_j$ for $c\ge2$. Then at every moment of the
run, $N_c\le A_c$ for all $c$, and
\[
A_c=\ell M\,(1+nM)^{\,c-1}.
\]
In particular every exponent occurring has absolute value at most $B=A_n=\ell M(1+nM)^{n-1}$,
of bit-length $\log_2 B=O\!\bigl(n(b+\log n)+\log\ell\bigr)$.
\end{proposition}

\begin{proof}
\emph{Solving the recursion.} For $c\ge2$, $A_c-A_{c-1}=nM\,A_{c-1}$, so $A_c=(1+nM)A_{c-1}$
with $A_1=\ell M$, giving $A_c=\ell M(1+nM)^{c-1}$. Taking $\log_2$,
\[
\log_2 B=\log_2\ell+b+(n-1)\log_2(1+n2^{b})=O\!\bigl(\log\ell+n(b+\log n)\bigr).
\]

\emph{The bound $N_c\le A_c$.} For each level $c$, define $T_c$ = the total level-$c$ mass
ever \emph{created} during the run, counting the initial input as creation at time $0$ and
counting each newly inserted syllable's $|e|$ once. Clearly $N_c\le T_c$ at every moment
(present mass never exceeds created mass, since masses are only created, transferred, or
destroyed). We bound $T_c$.

By Lemma~\ref{lem:level}, level-$c$ mass is created only:
\begin{itemize}
\item[(i)] initially, by the input: at most $\ell$ syllables of exponent $<M$, contributing
$\le \ell M$;
\item[(ii)] by an iteration acting on a level-$j$ syllable with $j<c$ (the only iterations
that can create level-$c$ mass, again by Lemma~\ref{lem:level}, since a level-$c$ or higher
iteration never creates level-$<c$, hence to create level $c$ the acting level must be $<c$;
moreover same-level ($j=c$) iterations only create level-$c$ residuals that conserve $N_c$, as
shown below, so they do not increase $T_c$ on net — we account for net creation).
\end{itemize}

Let us quantify (ii). Consider a single iteration acting on a level-$j$ syllable $(j,e)$ with
$j<c$. It is either a swap ($c'>j$ on the prefix side) or a power reduction:
\begin{itemize}
\item Power reduction at level $j$ inserts $|q|\le|e|$ copies of $w_{j,j}^{\pm1}$; each copy
contributes to $N_c$ at most the absolute value of the level-$c$ exponent of $w_{j,j}$, which
is $<M$, and there is at most one level-$c$ syllable per copy. So this iteration adds at most
$|e|\cdot M$ to $T_c$.
\item Swap at level $j<k$ inserts $\mathrm{Conj}(g_k^{f},g_j^{\sigma})$ consisting of $|f|$
copies of $w_{j,k}^{\pm1}$; each copy adds at most $M$ to $T_c$ (one level-$c$ syllable of
exponent $<M$). So this iteration adds at most $|f|\cdot M$ to $T_c$, where $f$ is the exponent
of the prefix syllable being consumed, $|f|\le N_k\le T_k$.
\end{itemize}

Now sum over the whole run. We introduce one accounting unit per syllable that is ever
\emph{removed} from the prefix $P$ or \emph{reduced} as a power, and bound the level-$c$ mass
each such unit can create. Concretely, every iteration that creates level-$c$ mass is, by
Lemma~\ref{lem:level}, either a swap acting at some level $j<c$ or a power reduction at some
level $j<c$; in each case it removes from circulation either the prefix syllable $g_k^{f}$
(a swap, consuming $|f|$ units of level-$k$ mass, $k>j$) or the syllable $g_j^{e}$ (a power
reduction, consuming the part $|e|-e_0\ge r_j$ of its level-$j$ mass). By the two bullet
points, each such removed unit of mass generates at most $M$ units of mass at any one fixed
higher level, hence at most $M$ units of level-$c$ mass.

It remains to bound the total mass removed at each level. Consider level $j$. Mass at level
$j$ is created only as in (i)--(ii), so the total level-$j$ mass ever in circulation is exactly
$T_j$. Each unit of it is removed at most once (once a syllable's mass is consumed by a swap or
reduction it is gone, replaced only by mass at strictly higher levels). Therefore the total
mass removed at level $j$ over the whole run is at most $T_j$. Summing the per-removed-unit
contribution $\le M$ to level $c$, over all levels $j<c$ from which level-$c$ mass can be fed,
and recalling that a single removed unit may, through a chain of swaps, feed at most $n$
distinct higher target levels (one per level above $j$), we obtain the uniform bound
\[
T_c\;\le\;\ell M\;+\;nM\sum_{j<c}T_j ,
\]
valid simultaneously for every $c$. This is exactly the defining recursion of $A_c$. By
induction on $c$ (the right-hand side involves only $T_j$ with $j<c$), $T_c\le A_c$, and
therefore $N_c\le T_c\le A_c=\ell M(1+nM)^{c-1}\le B$ for all $c$. Every individual exponent
has absolute value at most the mass of its level, hence at most $B$.
\end{proof}

\section{Number of operations, termination, and total running time}\label{sec:steps}

\begin{proposition}\label{prop:steps}
Algorithm~\ref{alg:collect} performs at most
\begin{equation}\label{eq:totalsteps}
T_{\mathrm{op}}\;\le\;\ell\,(1+nMB)^{\,n}\;=\;\ell\cdot 2^{\,O(n\log B)}\;=\;\ell\cdot 2^{\,O(n^2(b+\log(n\ell)))}
\end{equation}
elementary operations \textup{(}placements, merges, power reductions, swaps\textup{)}; in
particular it terminates. Each operation costs $O\!\bigl(n\,\mathsf M(\log B)\bigr)$ bit
operations, where $\mathsf M(t)=\widetilde O(t)$ is the cost of multiplying $t$-bit integers
and $\log B=O(n(b+\log(n\ell)))$. Hence the total running time is
\[
\ell\cdot 2^{\,O(n^2(b+\log(n\ell)))}=(nM\ell)^{O(n^2)},
\]
proving parts \textup{(1)}--\textup{(2)} of Theorem~\ref{thm:main}.
\end{proposition}

\begin{proof}
\emph{Counting operations by level.} Let $\mathcal O_c$ be the number of iterations whose
acting syllable has level $c$. We claim
\begin{equation}\label{eq:Oc}
\mathcal O_c\;\le\; 2\,T_c\;\le\;2A_c\;\le\;2B,
\end{equation}
where $T_c$ is, as in Proposition~\ref{prop:exp}, the total level-$c$ mass ever created.
Indeed, each level-$c$ iteration permanently removes work at level $c$: a placement moves a
level-$c$ syllable into $P$, where it stays for the rest of the run (placed syllables are only
ever removed again by a \emph{higher}-or-equal-level swap that consumes them as the prefix tail
— but such a swap is counted at \emph{its} acting level, not $c$); a power reduction at level
$c$ removes $\ge2$ units of level-$c$ mass; a merge removes one of two level-$c$ syllables; a
swap with acting level $c$ transfers exactly one unit of level-$c$ mass leftward toward
placement. Thus, charging each level-$c$ iteration to either a unit of level-$c$ mass it moves
toward $P$ or a unit it destroys, and since the total level-$c$ mass ever created is $T_c$,
the number of level-$c$ iterations is at most twice $T_c$ (the factor $2$ covers a swap moving
a unit left followed later by the placement of that unit). This gives \eqref{eq:Oc}.

\emph{Summing over levels.} By \eqref{eq:Oc} and Proposition~\ref{prop:exp},
\[
T_{\mathrm{op}}=\sum_{c=1}^{n}\mathcal O_c\;\le\;\sum_{c=1}^{n}2A_c
=2\sum_{c=1}^{n}\ell M(1+nM)^{c-1}
=2\ell M\,\frac{(1+nM)^{n}-1}{nM}
\;\le\;2\ell\,(1+nM)^{n}.
\]
Since $(1+nM)^n\le(1+nMB)^n$, this yields \eqref{eq:totalsteps}. Because $B$, hence each $A_c$,
is finite (Proposition~\ref{prop:exp}), $T_{\mathrm{op}}$ is finite: the algorithm performs
finitely many operations and therefore halts.

To see that $T_{\mathrm{op}}<\infty$ truly forces termination (and not an infinite run with
finitely many ``operations''), note every iteration of the while-loop either (a) discards a
zero syllable, (b) is one of the four elementary operations counted above, in which case it is
charged in $T_{\mathrm{op}}$, or (c) is a placement. Placements append to $P$ and are counted
in the level sum; discards strictly shorten $U$ and create nothing. As $T_{\mathrm{op}}$ bounds
all non-discard iterations and discards cannot occur infinitely often between two non-discard
iterations (each discard removes a syllable from $U$ that was inserted by an earlier counted
operation, and only finitely many syllables are ever inserted, namely at most
$\sum_c \mathcal O_c\cdot n\le nT_{\mathrm{op}}$ of them), the total number of while-loop
iterations is at most $(n+1)T_{\mathrm{op}}<\infty$. Hence the loop terminates.

\emph{Cost per operation.} An elementary operation manipulates the relator word being inserted
($\le n$ syllables) and a constant number of other syllables, performing $O(n)$ integer
additions/subtractions and at most one division-with-remainder (in a power reduction), all on
integers of absolute value $\le B$ by Proposition~\ref{prop:exp}, i.e.\ of bit-length $\le
\log_2 B+1=O(n(b+\log(n\ell)))$. Each such arithmetic step costs $O(\mathsf M(\log B))=
\widetilde O(\log B)$ bit operations with fast integer arithmetic (or $O((\log B)^2)$ with
schoolbook arithmetic). So the per-operation cost is $O(n\,\mathsf M(\log B))$.

\emph{Total.} Multiplying the operation count $(n+1)T_{\mathrm{op}}$ by the per-operation cost,
\[
(n+1)\,T_{\mathrm{op}}\cdot O\!\bigl(n\,\mathsf M(\log B)\bigr)
=\ell\cdot 2^{O(n\log B)}\cdot \widetilde O(n^2\log B)
=\ell\cdot 2^{O(n^2(b+\log(n\ell)))},
\]
using $\log B=O(n(b+\log(n\ell)))$ and absorbing the polynomial factor into the exponential.
This is \eqref{eq:upper}.
\end{proof}

\begin{remark}
For a \emph{fixed} presentation ($n=O(1)$, $b=O(1)$, hence $M=O(1)$), \eqref{eq:upper} reads
$\ell^{O(1)}$: collection is polynomial in the input length for each fixed polycyclic
presentation, recovering the classical fact. The blow-up is entirely in the dependence on $n$.
\end{remark}

\section{Lower bound: the exponential dependence on $n$ is genuine}\label{sec:lower}

\begin{proposition}\label{prop:lower}
For each $n\ge1$ let $G_n=\mathrm{UT}(n+1,\mathbb Z)$ be the group of upper unitriangular
$(n+1)\times(n+1)$ integer matrices. Then:
\begin{enumerate}
\item $G_n$ is polycyclic of Hirsch length $N=\binom{n+1}{2}=\tfrac{n(n+1)}{2}$, and admits a
consistent polycyclic presentation in which every relative order is $\infty$ and every relator
right-hand side is a normal word whose exponents lie in $\{-1,0,1\}$; thus $b=1$, $M=2$.
\item With $x_i=I+E_{i,i+1}$ \textup{(}elementary matrix, single $1$ in position $(i,i+1)$\textup{)}
for $i=1,\dots,n$, the input word $W_n=(x_1x_2\cdots x_n)^{n}$ has syllable length $\ell=n^2$,
and its normal form contains the generator $E_{1,n+1}=I+E_{1,n+1}^{\mathrm{mat}}$ with exponent
exactly $\binom{2n-1}{n}$.
\item $\binom{2n-1}{n}=\tfrac12\binom{2n}{n}\ge\frac{4^{n}}{4\sqrt n}$, so this exponent has
bit-length $\ge 2n-2-\tfrac12\log_2 n=\Theta(n)=\Theta(\sqrt N)$. Hence on $(G_n,W_n)$, with
$b=1$ and $\ell=n^2=\mathrm{poly}(n)$, the output normal form has a single exponent of
$\Theta(n)$ bits, so any correct algorithm — in particular Algorithm~\ref{alg:collect} — uses
time exponential in $n$ relative to the input size $\ell+b$.
\end{enumerate}
\end{proposition}

\begin{proof}
\emph{(1).} Order the $N$ elementary matrices $E_{i,j}$ ($1\le i<j\le n+1$) by increasing
weight $j-i$, breaking ties arbitrarily, and call them $h_1,\dots,h_N$ in this order. Each
$h_t=E_{i,j}$ has infinite order in $G_n$ (its powers are $I+m E_{i,j}^{\mathrm{mat}}$, all
distinct), so $r_t=\infty$ and there are no power relations. The conjugation relations are the
standard unitriangular commutator identities. Writing $[a,b]=a^{-1}b^{-1}ab$, in
$G_n$ one has
\[
[E_{i,j},\,E_{j,k}]=E_{i,k}\quad(i<j<k),\qquad
[E_{i,j},\,E_{k,l}]=I\quad\text{whenever }j\ne k\text{ and }i\ne l,
\]
and $E_{i,j}$ commutes with $E_{i,l}$ and with $E_{k,j}$ (matrices sharing a row or a column
index but not chaining). Consequently the conjugate $E_{p,q}^{\,E_{i,j}^{\pm1}}
=E_{i,j}^{\mp1}E_{p,q}E_{i,j}^{\pm1}$ equals $E_{p,q}$ multiplied by a (possibly empty) product
of elementary matrices $E_{a,b}$ each with exponent in $\{-1,0,1\}$: explicitly,
$E_{p,q}^{\,E_{i,j}^{\pm1}}=E_{p,q}$ unless $q=i$ (giving an extra factor $E_{p,j}^{\pm1}$) or
$p=j$ (giving an extra factor $E_{i,q}^{\mp1}$), in which cases the extra factor $E_{a,b}$ has
weight $b-a=(q-p)+(j-i)>q-p$, i.e.\ strictly larger weight than $E_{p,q}$. Thus every extra
factor is a higher-indexed generator in our weight ordering, so $E_{p,q}^{\,E_{i,j}^{\pm1}}$ is
already a normal word in the generators following $E_{p,q}$, with all exponents in
$\{-1,0,1\}$. Hence every conjugation relation $h_t^{\,h_s}$ ($s<t$) is a normal word in
$h_{t+1},\dots,h_N$ with exponents in $\{-1,0,1\}$, so $b=1$, $M=2$. Consistency holds because the matrix group
$G_n$ realizes the presentation faithfully and unitriangular matrices factor uniquely as
$\prod_{(i,j)}E_{i,j}^{a_{i,j}}$ in any fixed order of the $(i,j)$; thus the normal words are
in bijection with $G_n$, which is exactly consistency. The Hirsch length is the number of
generators, $N=\binom{n+1}{2}$.

\emph{(2).} Regard each word in the $x_i$ as its matrix. For unitriangular matrices, the
$(1,n+1)$ entry of a product $y_1y_2\cdots y_L$ of elementary factors $y_t=x_{m_t}$ equals the
number of strictly increasing index sequences $1\le p_1<p_2<\cdots<p_n\le L$ with
$y_{p_s}=x_s$ for $s=1,\dots,n$. (Proof: the $(1,n+1)$ entry counts directed paths
$1\to2\to\cdots\to(n+1)$ where the step $s\to s+1$ is contributed by some factor equal to
$x_s$, and the contributing factors must appear in increasing position because a later
matrix multiplies on the right; each factor $x_{m_t}=I+E_{m_t,m_t+1}^{\mathrm{mat}}$ can supply
at most the single step $m_t\to m_t+1$.) For $W_n=(x_1\cdots x_n)^n$ write the factors as
$y_{(\beta-1)n+m}=x_m$, $\beta=1,\dots,n$, $m=1,\dots,n$. The condition $y_{p_s}=x_s$ forces
$p_s=(\beta_s-1)n+s$ for some block index $\beta_s\in\{1,\dots,n\}$. The strict-increase
condition $p_1<\cdots<p_n$ reads $(\beta_1-1)n+1<(\beta_2-1)n+2<\cdots<(\beta_n-1)n+n$, which,
because the residues $1<2<\cdots<n$ are already increasing within a block, is equivalent to
$\beta_1\le\beta_2\le\cdots\le\beta_n$. The number of non-decreasing sequences $1\le
\beta_1\le\cdots\le\beta_n\le n$ equals $\binom{n+(n-1)}{n}=\binom{2n-1}{n}$. Thus the
$(1,n+1)$ entry of $W_n$ is $\binom{2n-1}{n}$. In the normal form, the exponent on $E_{1,n+1}$
equals this $(1,n+1)$ entry, because $E_{1,n+1}$ is the unique weight-$n$ generator and its
normal-form exponent is read directly off the matrix entry in position $(1,n+1)$.

\emph{(3).} We first show $\binom{2n}{n}\ge\frac{4^n}{2\sqrt n}$ for all $n\ge1$ by
induction. Base case $n=1$: $\binom{2}{1}=2=\frac{4}{2\sqrt1}$. Inductive step: assuming the
bound for $n$, and using $\binom{2n+2}{n+1}=\binom{2n}{n}\cdot\frac{(2n+1)(2n+2)}{(n+1)^2}
=\binom{2n}{n}\cdot\frac{2(2n+1)}{n+1}$, it suffices to verify
\[
\frac{4^n}{2\sqrt n}\cdot\frac{2(2n+1)}{n+1}\;\ge\;\frac{4^{n+1}}{2\sqrt{n+1}},
\]
which after cancelling $\frac{4^n}{2}$ and cross-multiplying is equivalent to
$(2n+1)\sqrt{n+1}\ge 2\sqrt n\,(n+1)$, i.e.\ (squaring, all quantities positive)
$(2n+1)^2(n+1)\ge 4n(n+1)^2$, i.e.\ $(2n+1)^2\ge 4n(n+1)$, i.e.\ $4n^2+4n+1\ge 4n^2+4n$, which
holds. This completes the induction. Hence
$\binom{2n-1}{n}=\tfrac12\binom{2n}{n}\ge\frac{4^n}{4\sqrt n}$, whose binary logarithm is
$\ge 2n-2-\tfrac12\log_2 n=\Theta(n)$. Since $N=\Theta(n^2)$ this is $\Theta(\sqrt N)$. The
input has size $\ell+b=n^2+1=\mathrm{poly}(n)$, yet the output contains an exponent of
$\Theta(n)$ bits, i.e.\ of magnitude $2^{\Theta(n)}$. Any algorithm correctly computing the
normal form must emit this exponent and therefore runs for at least $\Theta(n)$ steps merely to
write it; more to the point, no running-time bound of the form $\mathrm{poly}(\ell,n,b)$ with
$n$ appearing only polynomially can hold, since such a bound would be $\mathrm{poly}(n)$ here
whereas the required output length is $2^{\Theta(n)}$ for the magnitude (equivalently the
collection algorithm, which builds the exponent by additions of magnitude $\le 2^{\Theta(n)}$,
manipulates $\Theta(n)$-bit integers). This establishes the necessity of exponential dependence
on $n$.
\end{proof}

\begin{remark}[Sharpness]
Proposition~\ref{prop:lower} shows the blow-up is intrinsic to the \emph{output}: in
$\mathrm{UT}(n+1,\mathbb Z)$ a word of polynomial length over $\{\pm1\}$-exponent relators
produces normal-form exponents of magnitude $4^{\Theta(n)}$. This matches the upper bound
$B=\ell M(1+nM)^{n-1}=2^{\Theta(n)}$ of Proposition~\ref{prop:exp} when $\ell,M=O(1)$, up to
the constant in the exponent. Hence the worst-case complexity is genuinely
$\ell\cdot 2^{\Theta(\,\cdot\,)}$: polynomial in $\ell$ and $M=2^b$ for each fixed Hirsch
length $n$, and exponential in $n$.
\end{remark}

\section{Conclusion}

Combining Propositions~\ref{prop:exp}, \ref{prop:steps} and \ref{prop:lower} proves
Theorem~\ref{thm:main}. The worst-case complexity of the collection algorithm for computing
the normal form of a length-$\ell$ word in a polycyclic group of Hirsch length $n$, presented
by a consistent polycyclic presentation with exponent bit-bound $b$, is
\[
(nM\ell)^{O(n^2)}\;=\;\ell\cdot 2^{\,O(n^2(b+\log(n\ell)))},\qquad M=2^b,
\]
with all intermediate and final exponents of bit-length $O\!\bigl(n(b+\log(n\ell))\bigr)$; this
is polynomial in $\ell$ and in $M=2^b$ for every fixed Hirsch length $n$, and the exponential
dependence on $n$ is necessary, as witnessed by the unitriangular groups
$\mathrm{UT}(n+1,\mathbb Z)$.
\end{document}
